\section{Conclusions and Future Work}
In this paper we questioned the empirical realism of generative-agent-based modeling for imitating user behavior on social networks. First, we provided a formal framework for building realistic TWONs. Second, we instantiated this framework with the purpose of mimicking user behavior based on data from $\mathbb{X}$ in English and German. Third, we benchmarked the empirical realism of agents, imitating actual users.

\subsection{Key Recommendations}
Our empirical results provide several key recommendations for conducting social simulations based on generative-agent-based modeling:

First, simulation models should be validated with respect to their empirical realism before conducting simulations. Results of a simulation should always be put into perspective to the empirical realism of all components in the simulation.

Second, simulations should be performed in the same setting in which the simulation components were fitted and validated. The stark difference between English and German performance demonstrates this. The English models performed significantly closer to their real-world counterparts and produced more stable results. Changing the setting without retraining and validation can lead to unreliable outcomes of a simulation.

Third, fine-tuning of the simulation components is required to obtain sufficient levels of empirical realism. In-context prompting of LLMs, as is done in most of the related work \cite{larooij2025large}, specifically in areas like psychology, social sciences, or media studies is often not sufficient to guarantee a simulated behavior that is close to reality.

\subsection{Future Research}
This paper is a first step towards more robust empirical research designs and protocols for studying real-world social networks by simulating users with generative-agent-based models. Although we established general formal models and initial empirical insights, further research is required. The heterogeneous nature of social networks suggests that discourse patterns can vary significantly between different communities, each with its own linguistic norms, interaction styles, and argumentation preferences \cite{gnach2017social}. This diversity presents challenges both methodologically and theoretically. Methodologically, there is uncertainty about whether LLMs are adequate to capture these variations \cite{anderson2024homogenization} or if more specialized models tailored to specific communities are needed. Theoretically, this leads to a broader inquiry about whether social media discourse's inherent heterogeneity necessitates a community-specific modeling approach rather than universal models.
